\section{Formal Concept Analysis from scratch}
\label{sec:fca}

\fca{} was introduced by \citet{wille1982} as a way of ``restructuring lattice
theory'' --- taking an abstract branch of algebra and giving it a concrete
interpretation in terms of objects, attributes and concepts, so that lattice
diagrams become something a domain expert can read. The standard reference is
\citet{ganter1999}; \citet{carpineto2004} and \citet{priss2006} are gentler
surveys.

This section develops the theory on the running example of
Section~\ref{sec:running-example}. Everything is computed rather than
asserted: the numbers come from \file{fca_experiments.py}
(Appendix~\ref{sec:experiments}).

\subsection{The formal context}

\begin{definition}[Formal context]
A \emph{formal context} is a triple $\ctx = (G, M, I)$ where $G$ is a set of
\emph{objects}, $M$ is a set of \emph{attributes}, and
$I \subseteq G \times M$ is an \emph{incidence relation}. We write $g\,I\,m$
for ``object $g$ has attribute $m$''.
\end{definition}

The letters are from the German: \emph{Gegenst\"ande} (objects),
\emph{Merkmale} (attributes). A context is just a cross table, and our running
example is this one:

\begin{table}[h]
\centering
\small
\caption{The running example as a formal context $\ctx = (G, M, I)$: six
documents, five properties, $57\%$ of the cells filled.}
\label{tab:context}
\vspace{2mm}
\begin{tabular}{l ccccc}
\toprule
& \rotatebox{60}{\textsf{Flight}} & \rotatebox{60}{\textsf{Hotel}}
& \rotatebox{60}{\textsf{Japan}} & \rotatebox{60}{\textsf{Booked}}
& \rotatebox{60}{\textsf{Refundable}}\\
\midrule
\texttt{outbound.pdf}  & $\times$ &          & $\times$ & $\times$ & $\times$\\
\texttt{inbound.pdf}   & $\times$ &          & $\times$ & $\times$ &         \\
\texttt{kyoto-inn.pdf} &          & $\times$ & $\times$ & $\times$ & $\times$\\
\texttt{rail-pass.pdf} &          &          & $\times$ & $\times$ & $\times$\\
\texttt{paris-hop.pdf} & $\times$ &          &          & $\times$ &         \\
\texttt{visa-note.md}  &          &          & $\times$ &          &         \\
\bottomrule
\end{tabular}
\end{table}

Note what a context does \emph{not} contain. There is no hierarchy in it --- no
statement that \textsf{Flight} is a kind of anything, no statement that
\texttt{outbound.pdf} is more specific than \texttt{inbound.pdf}. The context
is raw co-occurrence. Everything hierarchical that follows will be
\emph{derived} from it. That is \fca{}'s central promise and, as we shall see,
also its central cost.

\subsection{The derivation operators}
\label{sec:derivation}

\begin{definition}[Derivation]
For a set of objects $A \subseteq G$ and a set of attributes $B \subseteq M$:
\[
  A' \;=\; \{\, m \in M \;:\; g\,I\,m \text{ for all } g \in A \,\},
  \qquad
  B' \;=\; \{\, g \in G \;:\; g\,I\,m \text{ for all } m \in B \,\}.
\]
$A'$ is ``the attributes shared by all of $A$''; $B'$ is ``the objects having
all of $B$''. Both are written with the same prime, and which one is meant is
determined by the argument.
\end{definition}

\begin{example}
In Table~\ref{tab:context}:
\begin{align*}
\{\texttt{outbound.pdf}, \texttt{inbound.pdf}\}' &=
  \{\textsf{Flight}, \textsf{Japan}, \textsf{Booked}\},\\
\{\textsf{Flight}, \textsf{Japan}\}' &=
  \{\texttt{outbound.pdf}, \texttt{inbound.pdf}\},\\
\{\textsf{Hotel}\}' &= \{\texttt{kyoto-inn.pdf}\}, \\
\varnothing' &= M \quad\text{(everything is shared by no objects at all).}
\end{align*}
\end{example}

The two maps form an (antitone) Galois connection in the sense of
Section~\ref{sec:closure}: enlarging $A$ can only shrink $A'$, and
$A \subseteq B' \iff B \subseteq A'$. Both composites are closure operators.
We write $A'' = (A')'$ and call it the closure of $A$.

\begin{pitfall}
The prime is not an inverse. $A''$ generally contains more than $A$: it is
$A$ plus every other object that happens to share all of $A$'s common
attributes. Try $A = \{\texttt{outbound.pdf}\}$: $A' =
\{\textsf{Flight},\textsf{Japan},\textsf{Booked},\textsf{Refundable}\}$ and
$A'' = \{\texttt{outbound.pdf}\}$ again --- but
$A = \{\texttt{visa-note.md}\}$ closes to $A'' =
\{\texttt{outbound.pdf}, \texttt{inbound.pdf}, \texttt{kyoto-inn.pdf},
\texttt{rail-pass.pdf}, \texttt{visa-note.md}\}$, because \texttt{visa-note}'s
only attribute is \textsf{Japan} and all those documents have it.
\end{pitfall}

\subsection{Formal concepts}

\begin{definition}[Formal concept]
A \emph{formal concept} of $\ctx$ is a pair $(A, B)$ with $A \subseteq G$,
$B \subseteq M$, $A' = B$ and $B' = A$. $A$ is the \emph{extent} (the things
falling under the concept), $B$ the \emph{intent} (the properties defining
it).
\end{definition}

The two conditions are a fixed-point demand, and they encode a very old
intuition about what a concept is: \emph{the extent and the intent must
determine each other exactly}. A concept is not any old pair of a set of
things and a set of properties. It is a pair in perfect balance: every object
in the extent has every attribute in the intent (that much is easy), and
crucially there is no object outside the extent that would also qualify, and
no attribute outside the intent that all the extent happens to share.

\begin{example}
$(\{\texttt{outbound.pdf}, \texttt{inbound.pdf}\},
\{\textsf{Flight}, \textsf{Japan}, \textsf{Booked}\})$ is a formal concept: the
two flights to Japan share exactly those three properties, and no other
document has all three. Informally, it is the concept \emph{Japanese flight
booking}, which nobody put in the table --- it emerged.

By contrast $(\{\texttt{outbound.pdf}\}, \{\textsf{Flight}\})$ is \emph{not} a
concept: \texttt{outbound.pdf} has more attributes than just \textsf{Flight},
so the pair is unbalanced.
\end{example}

\begin{proposition}
$(A,B)$ is a formal concept iff $A$ is a closed object set ($A''=A$) and
$B=A'$; equivalently iff $B$ is a closed attribute set ($B''=B$) and $A=B'$.
So concepts are in bijection with closed attribute sets, which is how
algorithms enumerate them.
\end{proposition}

\subsection{The concept lattice}
\label{sec:basic-theorem}

Order concepts by specificity: $(A_1, B_1) \leq (A_2, B_2)$ when
$A_1 \subseteq A_2$ --- equivalently, when $B_2 \subseteq B_1$. Fewer things,
more properties, further down. The set of all concepts under this order is
written $\Bcl(\ctx)$.

\begin{theorem}[Basic Theorem of \fca{}; \citealp{wille1982,ganter1999}]
\label{thm:basic}
$\Bcl(\ctx)$ is a complete lattice, in which
\[
  \bigwedge_{t \in T}(A_t, B_t) =
    \Bigl(\bigcap_{t} A_t,\ \bigl(\bigcup_{t} B_t\bigr)''\Bigr),
  \qquad
  \bigvee_{t \in T}(A_t, B_t) =
    \Bigl(\bigl(\bigcup_{t} A_t\bigr)'',\ \bigcap_{t} B_t\Bigr).
\]
Conversely, every complete lattice is isomorphic to the concept lattice of
some context. So ``complete lattice'' and ``concept lattice'' are the same
notion up to isomorphism.
\end{theorem}

Two things are worth extracting from that statement, one reassuring and one
expensive.

\emph{Reassuring:} meets are intersections of extents. ``The most specific
concept that is both a \textsf{Flight} and a \textsf{Japan} document'' always
exists, is unique, and is computed by intersecting. That is a strong
guarantee, and it is exactly what makes concept lattices navigable.

\emph{Expensive:} \emph{completeness} means \emph{every} subset has a meet,
including subsets nobody asked about. The lattice contains a concept for every
closed attribute set, whether or not any object realises it, whether or not
anybody would recognise it, and whether or not the co-occurrence that produced
it is a real regularity or a coincidence. Section~\ref{sec:explosion} measures
what that costs.

\subsection{Reading a line diagram}

\begin{figure}[t]
\centering
\dotfig[0.85]{lattice_travel}
\caption{The concept lattice of Table~\ref{tab:context}: 10 concepts, 13
covering edges, drawn with \fca{}'s standard \emph{reduced labelling}. Each
attribute and each object name appears exactly once. To read off whether an
object has an attribute, start at the object's node and walk \emph{upward}
along any lines: you have the attribute iff you can reach its label. The black
disc is a concept that introduces neither a new attribute nor a new object.
Diagram generated with Graphviz from the computed lattice.}
\label{fig:lattice}
\end{figure}

Figure~\ref{fig:lattice} shows the lattice. The convention deserves
explanation because it is initially surprising and then very economical.

In \emph{full labelling} you would write each concept's whole extent and whole
intent inside its node, which repeats every name many times. \emph{Reduced
labelling} writes each name once, at the concept where it is ``introduced'':

\begin{itemize}[leftmargin=1.4em]
\item attribute $m$ labels the concept $(m', m'')$ --- the \emph{largest}
  concept whose intent contains $m$;
\item object $g$ labels the concept $(g'', g')$ --- the \emph{smallest}
  concept whose extent contains $g$.
\end{itemize}

\noindent and then the reading rules recover everything:

\begin{itemize}[leftmargin=1.4em]
\item the \textbf{intent} of a concept is every attribute label reachable by
  walking \emph{up};
\item the \textbf{extent} is every object label reachable by walking
  \emph{down};
\item ``object $g$ has attribute $m$'' iff you can walk up from $g$'s label
  to $m$'s label.
\end{itemize}

\begin{tryit}
In Figure~\ref{fig:lattice}, find \texttt{rail-pass} and walk upward. You
should reach \textsf{Refundable}, \textsf{Japan} and \textsf{Booked}, but not
\textsf{Flight} or \textsf{Hotel}. Compare with Table~\ref{tab:context}. Then
find the concept where \textsf{Booked} sits and walk \emph{down}: you collect
all five booked documents --- and, satisfyingly, not \texttt{visa-note}.
\end{tryit}

\subsection{The full list, and what is in it}

Here is every concept of the running example, computed and sorted by extent
size. \texttt{Supp} is the extent size (in data-mining language, the support);
\texttt{stab} is Kuznetsov's \emph{stability} \citep{kuznetsov2007}, the
fraction of subsets of the extent that still close to the same intent --- a
measure of how much the concept depends on any particular object.

\begin{table}[h]
\centering\small
\caption{All 10 formal concepts of Table~\ref{tab:context}. Only 10 of the
$2^5 = 32$ conceivable attribute sets are closed.}
\label{tab:concepts}
\vspace{1mm}
\begin{tabular}{r l l r r}
\toprule
\# & \textbf{Intent} & \textbf{Extent} & supp & stab\\
\midrule
0 & $\varnothing$ & all six documents & 6 & 0.25\\
1 & \textsf{Booked} & all but \texttt{visa-note} & 5 & 0.38\\
2 & \textsf{Japan} & all but \texttt{paris-hop} & 5 & 0.50\\
3 & \textsf{Booked},\textsf{Japan} & \texttt{outbound},\texttt{inbound},\texttt{kyoto-inn},\texttt{rail-pass} & 4 & 0.38\\
4 & \textsf{Booked},\textsf{Flight} & \texttt{outbound},\texttt{inbound},\texttt{paris-hop} & 3 & 0.50\\
5 & \textsf{Booked},\textsf{Japan},\textsf{Refundable} & \texttt{outbound},\texttt{kyoto-inn},\texttt{rail-pass} & 3 & 0.62\\
6 & \textsf{Booked},\textsf{Flight},\textsf{Japan} & \texttt{outbound},\texttt{inbound} & 2 & 0.50\\
7 & \textsf{Booked},\textsf{Flight},\textsf{Japan},\textsf{Refundable} & \texttt{outbound} & 1 & 0.50\\
8 & \textsf{Booked},\textsf{Hotel},\textsf{Japan},\textsf{Refundable} & \texttt{kyoto-inn} & 1 & 0.50\\
9 & all five attributes & $\varnothing$ & 0 & 1.00\\
\bottomrule
\end{tabular}
\end{table}

Concepts 0 and 9 are the lattice's $\top$ and $\bot$: a complete lattice must
have them, so they are manufactured whether or not anything in the world
corresponds. Concept 9 is the concept of a document that is simultaneously a
flight and a hotel --- empty, and necessarily so. Concepts 3, 4, 5 and 6 are
the interesting ones: \emph{nobody wrote them down}, and each is a genuine
category (``booked Japan documents'', ``booked flights'', ``refundable Japan
bookings'', ``booked flights to Japan''). This is \fca{} doing what it is for.

\subsection{Implications}
\label{sec:implications}

\begin{definition}[Attribute implication]
For $B, C \subseteq M$, the implication $B \to C$ \emph{holds in} $\ctx$ if
every object having all attributes in $B$ also has all attributes in $C$;
equivalently, if $C \subseteq B''$.
\end{definition}

\begin{measured}
The single-premise implications of Table~\ref{tab:context}:
\[
\textsf{Flight} \to \textsf{Booked}, \qquad
\textsf{Refundable} \to \textsf{Booked}, \textsf{Japan}, \qquad
\textsf{Hotel} \to \textsf{Booked}, \textsf{Japan}, \textsf{Refundable}.
\]
\end{measured}

Read them aloud and the value is obvious: ``every flight here is booked'',
``everything refundable is a booked Japan document''. Read them again and the
danger is equally obvious: these are facts about \emph{this table}, not laws.
One unbooked flight destroys the first. Section~\ref{sec:learn-implications}
is about exactly this distinction, and it is the single most important thing
\odag{} must get right if it borrows anything from this part of \fca{}.

The set of all implications holding in a context is infinite (it is closed
under trivial consequences), but it has a canonical finite generating set:

\begin{theorem}[Duquenne--Guigues base; \citealp{guigues1986}]
Every formal context has a unique minimum-cardinality set of implications from
which all others follow, the \emph{stem base}. Its premises are the
\emph{pseudo-intents} of the context.
\end{theorem}

The stem base is a \emph{canonical} object: same context, same base. It is the
closest thing \fca{} has to \odag{}'s canonical form, and comparing the two
notions of canonicity (Section~\ref{sec:two-canonical-forms}) is one of the
more illuminating exercises in this paper.

\subsection{Computing: NextClosure and complexity}
\label{sec:algorithms}

The standard enumeration algorithm is \citeauthor{ganter1984}'s
\emph{NextClosure} \citeyearpar{ganter1984}. Fix an order on the attributes.
Every closed set has a successor in \emph{lectic} order (a variant of
alphabetical order on the characteristic bit vectors), computable directly.
Starting from $\varnothing''$ and repeatedly taking successors enumerates every
closed set, in order, using memory for one set at a time --- no lattice needs
to be held in RAM. The paper's implementation
(Appendix~\ref{sec:experiments}) is thirty lines and is checked against a
brute-force oracle on every example in this document.

Many faster algorithms exist --- Close-by-One and its descendants, Lindig's,
Titanic --- and they are compared empirically by \citet{kuznetsov2002} and
\citet{krajca2010}. None of them can escape the following.

\begin{theorem}[Hardness; \citealp{kuznetsov2001}]
The number of formal concepts of a context can be exponential in the size of
the context, and determining that number is $\#\mathrm{P}$-complete.
\end{theorem}

There is nothing pathological about the exponential case. It is reached by the
simplest imaginable table:

\begin{definition}[Contranominal scale]
$\mathbb{N}_n^c = (\{1,\dots,n\}, \{1,\dots,n\}, \neq)$: $n$ objects, $n$
attributes, and object $i$ has attribute $j$ exactly when $i \neq j$. Every
attribute subset is closed, so $|\Bcl| = 2^n$.
\end{definition}

\subsection{The explosion, measured}
\label{sec:explosion}

\begin{measured}
Contranominal scales, computed:

\smallskip
\centering\small
\begin{tabular}{rrrr}
\toprule
$n$ & concepts & \odag{} nodes & \odag{} edges\\
\midrule
 4 &      16 &  8 &  12\\
 8 &     256 & 16 &  56\\
10 &   1\,024 & 20 &  90\\
12 &   4\,096 & 24 & 132\\
\bottomrule
\end{tabular}
\end{measured}

\noindent
And on random data, where the effect is less extreme but still decisive.
Figure~\ref{fig:explosion} plots the mean number of concepts of a random
context with 30 objects and 12 attributes, against the number of crosses in
the table (which is what \odag{} would store as edges), as the density rises.

\begin{figure}[t]
\centering
\begin{tikzpicture}
\begin{axis}[
  width=0.78\textwidth, height=6.2cm,
  xlabel={density of the context (fraction of cells filled)},
  ylabel={count},
  xmin=0.05, xmax=0.65,
  legend pos=north west, legend cell align=left,
  grid=major, grid style={odgrey!20},
  every axis plot/.append style={very thick},
  tick label style={font=\small}, label style={font=\small},
  legend style={font=\small}]
\addplot[odrust, mark=*] table[x=density, y=concepts]
  {experiments/explosion.dat};
\addlegendentry{formal concepts (mean of 20)}
\addplot[odteal, mark=square*] table[x=density, y=edges]
  {experiments/explosion.dat};
\addlegendentry{crosses in the table $=$ \odag{} edges}
\end{axis}
\end{tikzpicture}
\caption{Random contexts, 30 objects $\times$ 12 attributes. The stored data
grows linearly with density; the concept lattice does not. At density $0.6$
the table has about 211 crosses and the lattice has about 550 concepts --- and
this is a \emph{small} context. Data from
\file{experiments/explosion.dat}, generated by the experiment script.}
\label{fig:explosion}
\end{figure}

This is not a criticism of \fca{}. Closing a context is an expensive
operation because \emph{the information is really there} --- those concepts
are genuinely supported by the data. It is, however, decisive for a system
that must hold a growing store in a browser tab, hash it after every edit, and
merge it with a stranger's copy. Section~\ref{sec:not-closed} returns to this.

\subsection{The parts of \fca{} beyond the basics}
\label{sec:fca-extensions}

\fca{} has grown a substantial toolbox in forty years. Six branches matter for
this paper; each reappears in Section~\ref{sec:learn} as a candidate for
adoption.

\paragraph{Many-valued contexts and conceptual scaling.} Real data has
numbers, not crosses. A \emph{many-valued context} $(G, M, W, I)$ assigns each
object a value in $W$ for each attribute; \emph{conceptual scaling} converts
it into an ordinary context by choosing, per attribute, a \emph{scale} that
says which comparisons matter \citep[ch.~1]{ganter1999}. The standard scales
are worth memorising because they are a checklist of ways to turn a value into
crosses:

\begin{center}\small
\begin{tabularx}{\textwidth}{@{}lXl@{}}
\toprule
\textbf{Scale} & \textbf{Turns a value into} & \textbf{Example}\\
\midrule
nominal      & one attribute per distinct value, mutually exclusive & colour\\
ordinal      & ``$\leq v$'' attributes, one per threshold & school grade\\
interordinal & both ``$\leq v$'' and ``$\geq v$'' & temperature ranges\\
biordinal    & two ordinal scales meeting in the middle & good/bad axes\\
contranominal& ``$\neq v$'' --- the exponential one & (avoid)\\
\bottomrule
\end{tabularx}
\end{center}

\paragraph{Pattern structures.} Scaling is lossy and arbitrary. Pattern
structures \citep{ganter2001} generalise \fca{} so that an object's
description need not be a set of attributes at all: it can be any element of a
meet-semilattice with a similarity operation $\sqcap$. The most-used instance
is the \emph{interval pattern structure}, where a description is an interval
and $\sqcap$ is the convex hull \citep{kaytoue2011}. Remember this; \odag{}'s
handling of typed values (\S\ref{sec:dimensions}) turns out to be an instance
of it that was arrived at independently.

\paragraph{Iceberg lattices and frequent closed itemsets.} Keep only concepts
whose extent exceeds a support threshold \citep{stumme2002}. The result is the
top of the lattice --- the ``iceberg'' --- and it is exactly the set of
frequent closed itemsets from association-rule mining \citep{pasquier1999}. A
principled way to be partial.

\paragraph{Stability.} A different filter: keep concepts that survive the
removal of objects \citep{kuznetsov2007}. Support asks ``how many?''; stability
asks ``how robust?''. The two disagree usefully, as
Table~\ref{tab:concepts} shows --- concept 5 has lower support than concept 3
but higher stability.

\paragraph{Boolean factor analysis.} \citet{belohlavek2010} proved that formal
concepts are exactly the optimal factors for exact Boolean matrix
decomposition: to write a binary matrix as a product of two binary matrices
with as few inner dimensions as possible, use concepts. This is the result
that connects \fca{} to the compression-based learning of
Section~\ref{sec:learn-mdl}: if concepts are factors, then choosing
\emph{which} concepts to keep is a model-selection problem, and MDL is the
principled answer \citep{miettinen2014,mdlfca}.

\paragraph{Attribute exploration.} An interactive protocol
\citep{ganter2016}: the system computes the next implication that its current
data does not refute, and asks a domain expert ``is this always true?''. The
expert either confirms --- and the implication joins the knowledge base --- or
supplies a counterexample, which is added as an object. Repeat. On
termination the base is complete for the domain, not merely for the sample.
Section~\ref{sec:learn-exploration} argues that this forty-year-old protocol
has just acquired an ideal new participant.

\paragraph{Further branches, in one line each.} \emph{Triadic} \fca{} adds a
third dimension --- objects $\times$ attributes $\times$ conditions
\citep{lehmann1995}. \emph{Fuzzy} \fca{} grades the incidence
\citep{belohlavek2004}. \emph{Relational} Concept Analysis handles links
between objects of different sorts by iterated scaling of relations
\citep{rouanehacene2013}. Applications range from software refactoring
\citep{snelting2000} to a broad survey of knowledge-processing uses
\citep{poelmans2013}.

\subsection{What to take away}

\begin{keyidea}
\fca{} is \textbf{derivation to completion}. Give it a binary table and it
returns every concept the table supports, organised into a complete lattice,
with a canonical implication base and a mature algorithmic toolbox. It answers
the question \emph{``what concepts does this data contain?''} exhaustively and
exactly.

Its costs are the direct price of that completeness: the answer can be
exponentially larger than the question; it describes one fixed table, so
adding a single object can restructure the whole lattice; and its concepts
have no names --- a concept \emph{is} its extent and intent, so when the data
changes, ``the same'' concept may not exist any more.
\end{keyidea}

Those three costs are precisely the three problems a live, shared, editable,
fingerprinted category store must not have. Section~\ref{sec:ontodag} shows
what a system looks like when it is designed around avoiding them.
